\documentclass[11pt,a4paper]{article}
\usepackage[T1]{fontenc}
\usepackage[utf8]{inputenc}
\usepackage[margin=2.5cm]{geometry}
\usepackage{lmodern}
\usepackage{microtype}
\usepackage{parskip}
\usepackage[table]{xcolor}
\usepackage{amsmath}
\usepackage[numbers]{natbib}
\usepackage[normalem]{ulem}
\usepackage{hyperref}
\usepackage{titlesec}
\hypersetup{colorlinks=true, linkcolor=blue, urlcolor=blue, citecolor=blue}

\definecolor{chgcolor}{RGB}{0,80,180}
\definecolor{rccolor}{RGB}{0,128,0}   % green for reviewer label

% ============================================================
% numbers.tex — gerado por make_all.py em 2026-10-05 23:17
% NÃO EDITAR À MÃO — editar make_all.py e regenerar
% ============================================================

% G1 — Qwen r=256 dose=0% N=3
\newcommand{\GoneKzero}{78}
\newcommand{\GoneMfive}{85.47}
\newcommand{\GoneSDfive}{0.75}
\newcommand{\GoneCIfiveLo}{83.60}
\newcommand{\GoneCIfiveHi}{87.30}
\newcommand{\GoneMten}{79.07}
\newcommand{\GoneSDten}{1.50}
\newcommand{\GoneCItenLo}{75.30}
\newcommand{\GoneCItenHi}{82.80}
\newcommand{\GoneSlope}{-1.28}
\newcommand{\GoneFourteenSlope}{-1.19}
\newcommand{\GoneSeedFifteenFive}{84.60}
\newcommand{\GoneSeedFifteenTen}{78.20}
\newcommand{\GoneSeedOneThreeSevenFive}{85.90}
\newcommand{\GoneSeedOneThreeSevenTen}{80.80}
\newcommand{\GoneSeedTwoFiveSixFive}{85.90}
\newcommand{\GoneSeedTwoFiveSixTen}{78.20}
% G1+G4 combinados N=6
\newcommand{\GoneFourMten}{79.72}
\newcommand{\GoneFourSDten}{1.73}
\newcommand{\GoneFourCItenLo}{77.90}
\newcommand{\GoneFourCItenHi}{81.50}
\newcommand{\GoneFourMfive}{85.68}
\newcommand{\GoneFourSDfive}{0.98}
% Testes seed-level Qwen r=16 vs r=256
\newcommand{\RsixteenMten}{97.40}
\newcommand{\RsixteenSDten}{0.00}
\newcommand{\DeltaRanks}{17.70}
\newcommand{\PermRanks}{0.0119}
\newcommand{\PermRanksBi}{0.0119}
\newcommand{\PermRanksCon}{0.0238}
\newcommand{\PermRanksNaloc}{84}
\newcommand{\HedgesGRanks}{10.94}
% G2 — Gemma3 r=10 dose-response N=3
\newcommand{\GtwoKzero}{44}
\newcommand{\GtwoMfiveZero}{83.33}
\newcommand{\GtwoMtenZero}{78.77}
\newcommand{\GtwoSDfiveZero}{1.33}
\newcommand{\GtwoSDtenZero}{1.27}
\newcommand{\GtwoMfiveTen}{89.37}
\newcommand{\GtwoMtenTen}{88.60}
\newcommand{\GtwoSDfiveTen}{1.33}
\newcommand{\GtwoSDtenTen}{0.00}
\newcommand{\GtwoDeltatenTen}{+9.8}
\newcommand{\GtwoCItenLoTen}{6.70}
\newcommand{\GtwoCItenHiTen}{13.00}
\newcommand{\GtwoPtenTen}{0.006}
\newcommand{\GtwoMfiveTfive}{90.90}
\newcommand{\GtwoMtenTfive}{90.13}
\newcommand{\GtwoSDfiveTfive}{0.00}
\newcommand{\GtwoSDtenTfive}{1.33}
\newcommand{\GtwoDeltatenTfive}{+11.4}
\newcommand{\GtwoCItenLoTfive}{5.80}
\newcommand{\GtwoCItenHiTfive}{17.00}
\newcommand{\GtwoPtenTfive}{0.013}
\newcommand{\GtwoMfiveFifty}{96.23}
\newcommand{\GtwoMtenFifty}{91.67}
\newcommand{\GtwoSDfiveFifty}{1.27}
\newcommand{\GtwoSDtenFifty}{1.33}
\newcommand{\GtwoDeltatenFifty}{+12.9}
\newcommand{\GtwoCItenLoFifty}{9.70}
\newcommand{\GtwoCItenHiFifty}{16.10}
\newcommand{\GtwoPtenFifty}{0.003}
\newcommand{\GtwoSlopeZero}{-0.91}
\newcommand{\GtwoSlopeFifty}{-0.91}
% G3 — NF4 vs bf16 confound
\newcommand{\GthreeDeltafiveSixteen}{+0.00}
\newcommand{\GthreeDeltatenSixteen}{+0.00}
\newcommand{\GthreeDeltafiveTwoFiveSix}{-0.20}
\newcommand{\GthreeDeltatenTwoFiveSix}{+0.00}
% G4 — variância de ordem
\newcommand{\GfourPorder}{0.700}
% Bloco 7 — dose-response Qwen N=5
\newcommand{\BsevenKzero}{78}
\newcommand{\BsevenMtenTen}{81.56}
\newcommand{\BsevenSDtenTen}{1.93}
\newcommand{\BsevenDeltatenTen}{+1.8}
\newcommand{\BsevenCItenLoTen}{-1.10}
\newcommand{\BsevenCItenHiTen}{4.70}
\newcommand{\BsevenPtenTen}{0.160}
\newcommand{\BsevenMtenTfive}{84.62}
\newcommand{\BsevenSDtenTfive}{2.21}
\newcommand{\BsevenDeltatenTfive}{+4.9}
\newcommand{\BsevenCItenLoTfive}{2.50}
\newcommand{\BsevenCItenHiTfive}{7.20}
\newcommand{\BsevenPtenTfive}{0.0044}
\newcommand{\BsevenMtenFifty}{88.98}
\newcommand{\BsevenSDtenFifty}{1.43}
\newcommand{\BsevenDeltatenFifty}{+9.2}
\newcommand{\BsevenCItenLoFifty}{6.90}
\newcommand{\BsevenCItenHiFifty}{11.60}
\newcommand{\BsevenPtenFifty}{0.001}
\newcommand{\BsevenSlopeZero}{-1.13}
\newcommand{\BsevenSlopeFifty}{-0.15}
\newcommand{\BsevenLRdelta}{+0.39}
% Gemma3 seed-level N=5+5
\newcommand{\GthreeKzero}{46}
\newcommand{\GthreeMtenRfour}{94.35}
\newcommand{\GthreeSDtenRfour}{1.19}
\newcommand{\GthreeMtenRsixteen}{56.52}
\newcommand{\GthreeSDtenRsixteen}{3.44}
\newcommand{\GthreeDeltaRanks}{37.80}
\newcommand{\GthreePermRanks}{0.0079}
\newcommand{\GthreeHedgesG}{13.29}
% G5 — Pi prospectivo
\newcommand{\GfiveRonetwoetightLRlow}{96.20}
\newcommand{\GfiveRonetwoetightLRhigh}{85.90}
\newcommand{\GfiveRthreytwoLRhigh}{89.70}



\titleformat{\section}{\normalsize\bfseries}{\thesection}{0.5em}{}

% ── Macros de estrutura IEEE Access ──────────────────────────────────────────
\newcommand{\rconcern}[3]{%
  % #1 = Reviewer #X, #2 = Concern #Y, #3 = short title
  \bigskip
  \noindent\textbf{\textcolor{rccolor}{#1, #2 (#3):}}%
  \par
}

\newcommand{\rquote}[1]{%
  \smallskip
  \begin{quote}\itshape #1\end{quote}%
}

\newcommand{\arhead}[1]{%
  \medskip\noindent\textbf{#1}\par\smallskip
}

\newcommand{\msloc}[1]{%
  \medskip\noindent\textit{Location in revised manuscript: #1}%
}

\begin{document}

% ── Cover page ────────────────────────────────────────────────────────────────
\noindent
Manuscript ID: \textbf{KNOSYS-D-26-21490}\\[2pt]
Original Article Title: \textit{Effective Training Pressure Gates Recursive
Knowledge Degradation in LLMs: A Multi-Axis Dose-Response Study}\\[2pt]
To: The Editor, \textit{Knowledge-Based Systems}\\[2pt]
Re: Response to Reviewers, Revision Round~1\\[2pt]
\today

\medskip\noindent\rule{\linewidth}{0.4pt}\medskip

\noindent
Dear Editor and Reviewers,

We thank you for your careful and constructive evaluation of our manuscript.
Your comments prompted substantive revisions to the experimental account,
statistical reporting, notation, and scope of our conclusions. We respond to
each comment below and indicate the corresponding changes in the manuscript.

The most consequential correction concerns the original synthetic-exposure
comparison. During revision, we identified that its baseline and intervention
conditions had been evaluated through different pipeline paths, including
differences in the construction of the baseline-correct question set
($K_0$). Consequently, we no longer rely on that comparison as evidence for
the previously reported intervention effect. The revised manuscript instead
presents paired dose--response experiments conducted under consistent
pipelines. On Qwen, reducing synthetic exposure attenuated the
Generation~5--10 retention decline at the 50\% reduction level. On Gemma~3,
the corresponding intervention improved Generation~10 retention but did not
attenuate the decline over that interval. We have revised the interpretation
to preserve this distinction rather than proposing a common mechanism that
the experiments do not establish.

We have also clarified the status of effective training pressure (ETP). It is
now presented as an operational framework for organizing experimental
conditions and their observed diagnostics, not as a validated universal
scalar predictor or an architecture-independent threshold. Additional
revisions distinguish controlled comparisons from descriptive contrasts,
report the available seed-level uncertainty, define the notation used in the
equations and figures, and state the limits of the factual-retention
evaluation.

The responses that follow identify the action taken for each concern and
direct the reader to the relevant part of the revised manuscript. Revised
manuscript passages quoted in this letter appear in
\textcolor{chgcolor}{blue}.

\bigskip
\noindent
Respectfully submitted on behalf of all authors,\\[6pt]
Julio Leite Azancort Neto\\
Corresponding Author

\medskip\noindent\rule{\linewidth}{0.4pt}\medskip

% ══════════════════════════════════════════════════════════════════════════════
\section*{Editor}

\rquote{Revise (including language editing). Provide point-by-point response
to each reviewer comment in the cover letter.}

\arhead{Author response:}
Thank you for this guidance. We have edited the manuscript throughout for
language and clarity, with particular attention to the description of
experimental conditions, the distinction between results and
interpretations, and the consistency of terminology across sections. We have
also moderated claims that extended beyond the tested architectures or
evaluation setting.

\arhead{Author action:}
This letter provides a separate response to every concern raised by
Reviewers~1 and~2. Each response describes the revision or, where a
requested conclusion cannot be established from the available experiments,
states that limitation explicitly. Locations in the revised manuscript are
provided with the individual responses below.

% ══════════════════════════════════════════════════════════════════════════════

\section*{Reviewer 1}

\rconcern{Reviewer \#1}{Concern \#1}{Prospective calibration of ETP}

\rquote{How can effective training pressure be quantitatively calibrated and
prospectively tested on an unseen backbone or dataset?}

\arhead{Author response:}
We agree that a quantitative framework should not be presented as
prospectively predictive without validation on held-out settings. The revised
manuscript therefore distinguishes variables specified before training
(nominal adapter rank, learning rate, planned training steps, and synthetic
exposure) from quantities measured afterward (effective rank and weight
drift). Retention, response length, and SDI-3 are evaluated as outcomes or
diagnostics, not used to define a prospective threshold.

Within the tested protocols, we can locate changes in retention trajectories
between adjacent configurations. Those intervals do not establish a universal
numerical ETP threshold: the attempted cross-architecture normalisations did
not align the observed boundaries. A prospective test on an unseen backbone
or dataset would require a calibration rule and outcome criterion specified
in advance, followed by evaluation on held-out conditions without refitting
the rule to their results. We have added this as a proposed validation
protocol, not as an experiment already completed.

The exploratory G5 experiment provides a narrower check: two of three
pre-specified, single-seed configurations followed the expected ordering. No
threshold was pre-specified for G5, and it did not evaluate an unseen
backbone or dataset. We therefore describe ETP as an operational
experimental framework rather than a validated, transferable predictor.

\arhead{Author action:}
We revised the ETP definition, qualified the G5 interpretation, and added
the requirements for future prospective calibration.

\msloc{Section~3.5 (ETP definition); Section~4.4 (G5); Section~6
(prospective validation and limitations).}

\rconcern{Reviewer \#1}{Concern \#2}{Separating pressure effects from confounders}

\rquote{How are quantization, optimizer state, data order, and model-specific
dynamics separated from the claimed pressure effects?}

\arhead{Author response:}
The revised manuscript identifies what each control can establish and what
remains confounded. For quantization, G3 compares Qwen at \(r=16\) and
\(r=256\) under the same bf16 protocol and three paired seeds. The
per-seed Generation~10 differences are \(+15\), \(+13\), and \(+15\)
questions out of 78, a mean difference of \(+18.4\) percentage points.
The exact paired sign-flip test gives \(p=0.25\). This directionally
consistent, small-sample result supports the rank contrast under bf16 but
does not establish a statistically conclusive or architecture-independent
effect.

By contrast, the \(17.7\)-percentage-point Qwen difference obtained by
combining conditions from different protocols is descriptive: rank,
quantization, and aspects of the data-order protocol differ across the
groups. We do not use that comparison as a rank-only hypothesis test.
Comparing NF4 and bf16 provides an additional implementation check, not
a complete causal account of the observed trajectories.

A fresh optimizer is instantiated for each generation, so optimizer state
is not carried across generations under the reported protocol. In G4,
both the random seed and data order changed relative to G1. The
Generation~10 comparison (\(N=3+3\), permutation \(p=0.70\)) consequently
does not isolate an order effect; nor does the nonsignificant result prove
that order has no effect.

Model-specific dynamics likewise remain relevant. For Gemma~3, the
\(r=4\) versus \(r=16\) contrast involves five runs per condition. An
independent-seed exact permutation test yields \(p=0.0079\), whereas
a paired sign-flip analysis yields \(p=0.0625\). We report both results
with their differing seed-assignment assumptions rather than treating
either test as universally applicable.

\arhead{Author action:}
We separated design variables from post-training diagnostics, documented
the controls and their limitations, and distinguished the paired G3
comparison from the descriptive cross-protocol Qwen contrast.

\msloc{Sections~3.5 and~3.9 (variables and controls); Section~4.2
(Qwen comparisons); Section~4.3.1 (Gemma~3); Section~4.6 (G3 and G4).}

\rconcern{Reviewer \#1}{Concern \#3}{External validity of the evaluation}

\rquote{Do the conclusions hold for longer answers, multi-hop questions,
other domains, and larger factual evaluation sets?}

\arhead{Author response:}
Our results do not establish generalization to those settings. The
evaluation concerns short-answer factual questions. It uses a held-out set
of 200 questions disjoint from the 2,000-question synthetic training
stream; retention is calculated over questions answered correctly at
generation~0 (\(|K_0|=78\) for Qwen under the reported protocol).

The primary scorer permits bidirectional substring matches to accommodate
short-answer variants. Because longer generated answers could benefit from
that permissive criterion, we also report a strict exact-match
re-evaluation of saved predictions. Absolute scores are lower under the
strict scorer, while the reported qualitative trends for the evaluated
conditions are retained. This sensitivity check addresses one scoring
concern; it is not a test of long-form generation or multi-hop reasoning.

We have corrected the discussion of evaluation-set size. Expanding a
held-out, disjoint evaluation set would not, by itself, change the
composition of the training stream. Larger factual test sets would improve
the assessment of sampling variability and generality, while other
domains and answer formats require separate experiments.

\arhead{Author action:}
We clarified the scorer and the training--evaluation separation, added
the strict-scoring comparison, and explicitly limited the conclusions to
the tested short-answer factual setting.

\msloc{Section~3.1 (evaluation protocol); supplementary Table~S2
(strict scoring); Section~6 (scope and limitations).}

\rconcern{Reviewer \#1}{Concern \#4}{Exposure reduction and model-dependent responses}

\rquote{Why does the five-percent exposure reduction stabilize Qwen but not
Gemma~3, and does this generalize across ranks and seeds?}

\arhead{Author response:}
This question exposed a substantive problem with the original comparison.
C1 and C5 followed different evaluation pipelines, including differences
in \(K_0\) construction and data loading. The previously reported
\(9.0\)- and \(9.4\)-percentage-point intervention gains cannot therefore
serve as evidence of stabilization. We also clarified the dose terminology:
C5 removed 15\% of training \emph{examples}, corresponding to
approximately 5\% of \emph{tokens}. Under a matched pipeline with three
paired seeds, the estimated effect of that smaller intervention was
\(+1.9\) percentage points and was inconclusive.

The revised Qwen evidence comes from the B7 dose--response experiment:
0\%, 10\%, 25\%, and 50\% example removal at \(r=256\), with five paired
seeds per dose. Relative to no removal, Generation~10 retention increased
by \(+1.8\) percentage points at 10\% removal (\(p=0.160\)),
\(+4.9\) points at 25\% (\(p=0.005\), 95\% CI \([2.5,7.2]\)),
and \(+9.2\) points at 50\% (\(p<0.001\), 95\% CI \([6.9,11.6]\)).
At 50\% removal, the Generation~5--10 retention slope was \(-0.15\)
percentage points per generation, compared with \(-1.13\) without
removal. Thus, attenuation of the later decline is supported for this
Qwen configuration and protocol; it is not inferred from the invalid
C1/C5 comparison.

Gemma~3 showed a different response. In G2, 50\% removal at \(r=10\)
increased Generation~10 retention by \(7.7\) percentage points across
three paired seeds, but the Generation~5--10 slope was approximately
\(-0.91\) percentage points per generation in both arms. The intervention
improved endpoint retention without demonstrably slowing later loss.
Neither experiment identifies the mechanism or establishes that the
effect generalizes across untested ranks. Similar point estimates for
50\% exposure removal and a halved learning rate are not presented as
formal equivalence.

\arhead{Author action:}
We disclosed the pipeline inconsistency, withdrew the original C1/C5
inference, replaced it with the paired B7 dose--response analysis, and
reported the distinct Gemma~3 outcome without attributing both patterns
to the same mechanism.

\msloc{Section~4.5 (pipeline correction, B7, G2, and learning-rate
comparison); Section~6 (mechanistic and generalization limits).}

\rconcern{Reviewer \#1}{Concern \#5}{Scope of the reproducibility package}

\rquote{What complete code, checkpoints, prompts, seeds, and raw outputs
will support independent reproduction?}

\arhead{Author response:}
We agree that aggregate tables alone are insufficient and that the scope
of the public materials must be described precisely. The revised
reproducibility description identifies the model checkpoints, question
partitions and baseline-correct \(K_0\) sets, generation and evaluation
prompts, decoding settings, training configurations, seeds, and the
per-seed results used in the reported analyses. It also directs readers
to the associated archive at
\url{https://doi.org/10.5281/zenodo.23146342}.

Importantly, we distinguish archived analysis materials from artifacts
that are not publicly deposited. The described release includes
consolidated results and replication scripts, but raw token-level
predictions, trained adapter checkpoints, and base-model weights are not
included. We therefore do not claim that the archive alone permits a
complete token-level audit or byte-for-byte reconstruction of every
training run. Scripts introduced or revised after release
\texttt{v1.0.3} must be identified by the version in which they are
actually deposited; they should not be attributed to the earlier release.

\arhead{Author action:}
We expanded the reproducibility and data-availability description to
state what the public package contains, what it does not contain, and
which archive version supports each reproduction claim.

\msloc{Appendix~A.1 (reproducibility and data availability).}


\section*{Reviewer 2}

% ─────────────────────────────────────────────────────────────────────────────
\rconcern{Reviewer \#2}{Concern \#1}{Abbreviations and notation}

\rquote{ETP appears in Figure~1 without explicit definition. MTLD is not
expanded. KL, JS, and coverage appear in Figure~1 without definition. bf16,
SVD, TCE, QLoRA, and LoRA lack consistently expanded forms. ``Homeo.'' is
used in Table~2 without definition. $K_0$ is used both as a set and as a
retention metric. $T$, $D_{1,T}$, $D_{10}$, $\Delta I$, and related symbols
in Eq.~6 are not fully defined.}

\arhead{Author response:}
Each notation issue has been corrected. The revised manuscript does not use
``ETP'' as a standalone acronym; the term is introduced in full at l.~72 and
used in full throughout. Figure~1 no longer contains the abbreviation in
isolation. MTLD now reads \textcolor{chgcolor}{``Measure of Textual Lexical Diversity (MTLD)''}
at first occurrence. KL, JS, and coverage were conceptual labels for
diagnostics never computed or reported; they have been removed from Figure~1
entirely, and the revised caption describes only quantities that are measured.
All acronyms now have explicit parenthetical expansions at first use:
\textcolor{chgcolor}{``full fine-tuning (FFT)''} at l.~75,
\textcolor{chgcolor}{``singular value decomposition (SVD)''} at l.~332,
\textcolor{chgcolor}{``Quantized Low-Rank Adaptation (QLoRA)''} at l.~333. TCE (Truncated
Cross-Entropy) is cited in the Related Work section as a mitigation method
from prior literature; it does not appear in the Methods or Results.
``Homeo.'' has been expanded to \textcolor{chgcolor}{``Homeostatic''} in all four occurrences
in Table~2.

A dedicated notation table (Table~1) distinguishes the two uses of $K_0$:
$K_0$ denotes the \emph{set} of held-out questions answered correctly at
generation~0; $|K_0|$ denotes its cardinality; and $R(t)=|K_t|/|K_0|$
denotes the retention ratio. All symbols in Eq.~6 (SDI-3) are defined in
Table~1 and in the inline text preceding the equation: $T$ = evaluation
horizon in generations; $D_{1,t}$ = Distinct-1 unigram diversity at
generation~$t$; $\Delta I = \mathrm{Instability}(T) - \mathrm{Instability}(1)$,
where $\mathrm{Instability}(t)$ is the fraction of training items whose
normalized response differs between generation~$t{-}1$ and generation~$t$;
$\mathrm{MeanLen}_t$ = mean response length in words at generation~$t$.

\arhead{Author action:}
We removed ``ETP'' from Figure~1, expanded MTLD and all acronyms at first use,
removed KL/JS/coverage from Figure~1, expanded ``Homeo.'' throughout Table~2,
added Table~1 (tab:notation) with all symbol definitions, and added inline
definitions for all Eq.~6 quantities in the preceding prose.

\msloc{Abstract l.~72; Figure~1 caption ll.~692--698; Table~1 (tab:notation)
ll.~438--445; Equations~1, 3, 5, 6; l.~875 (MTLD); ll.~75, 332, 333
(acronyms); Table~2 regime column.}

% ─────────────────────────────────────────────────────────────────────────────
\rconcern{Reviewer \#2}{Concern \#2}{Equations and tables}

\rquote{Equations~1--6 are severely malformed in the PDF. Eq.~1 uses
``Retention.t /'' instead of a clear function notation. Eq.~2 for effective
rank is missing the negative sign and summation. Eq.~6 for SDI-3 has unclear
signs and terms. ``$5\times10^6$'' in Table~7 appears to be a rendering
error.}

\arhead{Author response:}
The malformed rendering was a font-encoding artefact: Unicode characters in
mathematical mode were not mapped correctly to the T1 font metric, producing
garbled output while the source was syntactically correct. The revised
submission uses the \texttt{lmodern} package with T1 encoding throughout,
which eliminates all rendering artefacts.

Eq.~1 now reads correctly:
\[
  R(t) = \frac{|\{q \in K_0 : q \text{ correct at gen.\ } t\}|}{|K_0|}.
\]

Eq.~2 now includes the negative sign and the explicit normalisation:
\[
  \mathrm{erank}(BA) = \exp\!\left(-\sum_i \hat{\sigma}_i \log \hat{\sigma}_i\right),
\]
where $\hat{\sigma}_i = \sigma_i / \sum_j \sigma_j$. This is the standard
effective-rank definition of \citet{roy2007}.

Eq.~6 now reads:
\[
  \mathrm{SDI\text{-}3} =
  \log\frac{\mathrm{MeanLen}_T}{\mathrm{MeanLen}_0}
  + \log\frac{\mathrm{D1}_0}{\mathrm{D1}_T}
  + \Delta I.
\]
A positive SDI-3 score indicates cumulative distributional deterioration.
The value in Table~7 is the learning rate $5\times10^{-6}$; the exponent
$-6$ was dropped by the same encoding issue and is now rendered correctly.

\arhead{Author action:}
We added \texttt{lmodern} with T1 encoding to the preamble, rewrote Equations
1, 2, and 6 with explicit notation consistent with Table~1, and verified that
no other table or equation contains a similar rendering failure in the revised
submission.

\msloc{Preamble (lmodern + T1); Eq.~1 l.~595; Eq.~2 ll.~798--803; Eq.~6
ll.~881--886; Table~7 (tab:ranklr) l.~1473.}

% ─────────────────────────────────────────────────────────────────────────────
\rconcern{Reviewer \#2}{Concern \#3}{Technical and data inconsistencies}

\rquote{(a) Abstract and Section~4.2.3 state the Qwen boundary is between
effective ranks 50--88; Section~5.1 states Qwen remains outside the
degradative regime ``up to effective rank~150,'' contradicting Table~2.
(b)~$r=128$ is classified as Bounded in Table~2 but called
``above-threshold'' in Section~4.4.3. (c)~$r=256$ is repeatedly called a
``boundary configuration'' but Table~2 classifies it as Degradative.
(d)~Section~5.2 states ``$r\geq16$ on Gemma~3'' but Table~4 shows
$r\in\{10,12,14\}$ already degradative at Gen~5.
(e)~Abstract overclaims ``three to five independent seeds per condition''
when many conditions use a single seed. (f)~FFT on Gemma~3 is ``confirmed''
from a single-seed non-monotonic result. (g)~C5 removal percentage
discrepancy: 15\% examples vs.\ ${\sim}5\%$ tokens.}

\arhead{Author response:}
Each inconsistency has been corrected. The phrase ``up to effective rank~150''
has been removed; all threshold language now uses \textcolor{chgcolor}{``an abrupt,
threshold-like transition within the tested grid''}, reflecting empirical
localisation between adjacent tested configurations.

``Above-threshold configurations'' in Section~4.4.3 now refers exclusively
to $r=256$ (Degradative). The $r=128$ configuration (Bounded) is described
as a \textcolor{chgcolor}{``transition-zone configuration''} where distributional signatures
begin to emerge without producing the full degradative trajectory. The five
occurrences of ``boundary configuration'' referring to $r=256$ have been
replaced with \textcolor{chgcolor}{``lowest-pressure degradative configuration tested''}.

Table~4 reports retention of 78.3\,\%, 73.9\,\%, and 69.6\,\% for $r=10$,
12, and 14 at Gen~5, all Degradative. The text now reads
\textcolor{chgcolor}{``the Gemma~3 regime boundary lies between effective ranks 3 and 6, below
$r=10$''}; ``$r\geq16$'' has been replaced with \textcolor{chgcolor}{``$r\geq10$ on Gemma~3''}.

The abstract now reads \textcolor{chgcolor}{``three to five independent seeds for headline
conditions (single seed for intermediate ablations)''}, with a dedicated $N$
column added to all result tables. For FFT on Gemma~3, ``confirms'' has been
replaced with \textcolor{chgcolor}{``is consistent with''} and the non-monotonic pattern
(78.3\%, 54.3\%, 63.0\%) is acknowledged explicitly. The 15\%/5\%
discrepancy is clarified: C5 removed 15\% of training \emph{examples},
reducing the token count by ${\sim}5\%$ because responses vary in length.
The more substantive resolution is that the C1/C5 pipeline inconsistency
renders the original comparison uninformative; the B7 dose--response replaces
it.

\arhead{Author action:}
We removed ``up to effective rank~150,'' replaced all five occurrences of
``boundary configuration'' for $r=256$, corrected the $r=128$ label,
updated the Gemma~3 boundary description, corrected the abstract seed-count
claim, softened the FFT Gemma~3 conclusion, and clarified the example-vs.-token
distinction for C5.

\msloc{Abstract ll.~69, 72; Contribution~\#1 l.~95; Section~4.4.3
ll.~1061, 1570; ll.~1737, 1915, 2063, 2069, 2159; Section~5.2; Table~4
ll.~1265, 1277--1280; ll.~1147--1149 (FFT Gemma~3); Section~3.7 (C5
doses); Section~4.5 ll.~1708--1713.}

% ─────────────────────────────────────────────────────────────────────────────
\rconcern{Reviewer \#2}{Concern \#4}{Methodology}

\rquote{No formal statistical tests or confidence intervals are reported. The
ETP framework is not quantified; it is a conceptual label. Cross-architecture
normalisation failed with no alternative proposed. ``Effective training
pressure'' is not a new quantitative measure. Gemma~4 E2B lacks a citation.
The ``sharp transition'' claim is weakened by sparse rank grids.}

\arhead{Author response:}
The revised manuscript adds seed-level exact permutation tests for all
comparisons where independent training runs are available. For Qwen, the
17.7\,pp arithmetic gap between $r=16$ and $r=256$ is reported as a
descriptive observation because the two groups are not exchangeable under a
rank-only null. The cleanest within-protocol comparison is G3 paired design:
mean $+18.4$\,pp, sign-flip $p=0.25$ ($N=3$). For Gemma~3, bilateral
permutation $p=0.0079$ (independent seeds) or sign-flip $p=0.0625$ (paired).
For dose--response, paired $t$-tests with 95\,\% CIs are reported per dose
level.

ETP is presented as an \emph{operational experimental framework}, not a
validated scalar predictor. We distinguish ETP-design variables (rank, LR,
exposure fraction, steps) from ETP-observed diagnostics (effective rank,
$d(t)$, response length, SDI-3). The $\mathrm{erank}\times\mathrm{LR}$ index
is labelled as exploratory throughout. The failed normalisations in Table~6
are retained as evidence of incommensurability; architecture-specific
calibration is explicitly deferred to future work. The contribution is the
factorial design, the dose--response evidence, the pre-registered exposure
result, and the identification of architecture-specific regime boundaries.

The Gemma~4 E2B backbone is now cited as \citet{gemma4_2026} (Gemma Team,
2026, arXiv:2607.02770) at every introduction point. ``Sharp transition'' has
been replaced throughout with \textcolor{chgcolor}{``abrupt, threshold-like transition within
the tested grid''}.

\arhead{Author action:}
We added permutation tests for Qwen G3 and Gemma~3, paired $t$-tests with
CIs for all B7 and G2 dose levels, revised Section~3.5 to formalise the
ETP-design/ETP-observed distinction, added the Gemma~4 citation to
Sections~3.3 and~4.3, and replaced ``sharp transition'' throughout.

\msloc{Section~4.2 (G3 paired contrast, $p=0.25$); Section~4.3.1 (Gemma~3,
$p=0.0079$/$p=0.0625$); Section~4.5 (dose--response, paired $t$-tests);
Section~3.5 (ETP framework); Table~1 (tab:notation); Section~6
(cross-architecture limitation); Section~3.3 (Gemma~4 citation);
Abstract l.~72.}

% ─────────────────────────────────────────────────────────────────────────────
\rconcern{Reviewer \#2}{Concern \#5}{Format and citations}

\rquote{References [12] and [33] share arXiv:2408.14572. Reference [32] lists
year 2024 but the arXiv ID is from December 2025. No citation for Gemma~4
E2B. Keisha et al.\ is a preprint with no venue. Reference formatting is
inconsistent.}

\arhead{Author response:}
The arXiv number 2408.14572 now appears exactly once, assigned to
\citet{fawi2024curlora} (\textit{CURLoRA}, Fawi, 2024); the duplicate entry
has been removed. The entry \texttt{ding2024ranktrade} (Rathore et al.) now
carries \texttt{year = \{2025\}}, consistent with its December~2025
arXiv submission date and its appearance in AACL-IJCNLP~2025. The Gemma~4
E2B backbone is now cited as \citet{gemma4_2026} at every introduction and
description point. \citet{keisha2025} is a non-peer-reviewed preprint
(arXiv:2509.04796), now declared at first citation:
\textcolor{chgcolor}{``\citet{keisha2025} (non-peer-reviewed preprint) report...''}; no
quantitative claim depends exclusively on it; \citet{shumailov2024}
(published in \textit{Nature}) is added as the primary citation for the
model-collapse phenomenon.

\arhead{Author action:}
We audited and harmonised the full reference list: published works include
venue, volume, pages, and DOI; arXiv-only works are formatted as
\texttt{arXiv preprint arXiv:XXXX.XXXXX}. Entries
\texttt{ding2024ranktrade}, \texttt{fawi2024curlora}, and
\texttt{zibakhsh2024} have been completed. The full revised bibliography
is provided as \texttt{cas-refs.bib} in the submission package.

\msloc{References (cas-refs.bib): entries \texttt{fawi2024curlora},
\texttt{ding2024ranktrade}, \texttt{gemma4\_2026}, \texttt{keisha2025},
all entries harmonised.}

\end{document}
